\documentclass[a4paper,12pt]{article}
\usepackage[T2A]{fontenc}
\usepackage[utf8]{inputenc}
\usepackage[russian]{babel}
\usepackage{amsmath,amssymb,mathtools}
\usepackage{helvet}
\renewcommand{\familydefault}{\sfdefault}
\usepackage{icomma}
\usepackage[a4paper,top=23mm,bottom=23mm,left=23mm,right=23mm,headheight=16pt]{geometry}
\usepackage{microtype}
\usepackage{tikz}
\usetikzlibrary{arrows.meta,positioning,shapes.geometric,calc}
\usepackage{tabularx,booktabs}
\usepackage{enumitem}
\usepackage{parskip}
\usepackage{fancyhdr}
\usepackage{setspace}
\usepackage{xcolor}
\usepackage{hyperref}
\definecolor{accent}{HTML}{008080}
\definecolor{darkaccent}{HTML}{005B60}
\definecolor{textgray}{HTML}{353535}
\definecolor{softgray}{HTML}{6C7579}
\hypersetup{hidelinks}
\setstretch{1.20}
\setlength{\parindent}{0pt}
\setlength{\parskip}{6pt}
\setlist[enumerate]{leftmargin=1.5em,itemsep=7pt,topsep=5pt}
\setlist[itemize]{leftmargin=1.55em,itemsep=3pt,topsep=4pt}
\pagestyle{fancy}
\fancyhf{}
\fancyhead[L]{\footnotesize\color{softgray}Основы комбинаторики}
\fancyhead[R]{\footnotesize\color{softgray}Екатерина Гнедкова}
\fancyfoot[C]{\small\color{softgray}\thepage}
\renewcommand{\headrulewidth}{0pt}
\newcommand{\sectionline}[1]{\vspace{0.6em}\noindent{\large\bfseries\color{darkaccent}#1}\par\vspace{0.12em}}
\newcommand{\subline}[1]{\noindent{\bfseries\color{darkaccent}#1}\par}
\newcommand{\keyidea}[1]{{\color{darkaccent}\bfseries Ключевая идея.} #1}
\newcommand{\example}[1]{{\bfseries Пример.} #1}
\newcommand{\task}[1]{\item #1}
\begin{document}
\thispagestyle{empty}
\begin{center}
\vspace*{15mm}
{\large\color{darkaccent} МАТЕМАТИКА}\par
\vspace{19mm}
{\LARGE\color{textgray}Чек-лист}\par
\vspace{11mm}
{\Large\color{accent} Основы комбинаторики}\par
\vspace{3mm}
{\large\color{accent} Метод позиций и правило умножения}\par
\vspace{14mm}
{\normalsize Подготовка к ЕГЭ по профильной математике}\par
\vspace{30mm}
{\normalsize Екатерина Гнедкова}\par
\vspace{10mm}
{\normalsize 28 октября 2026 г.}\par
\vfill
{\small Лёвин Артём Александрович}\par
{\small эксклюзивно для Екатерины Гнедковой}\par
\vspace{13mm}
\begin{tikzpicture}[x=1cm,y=1cm]
  \draw[accent,line width=.9pt] (0,0) .. controls (1.1,.55) and (2.3,-.5) .. (3.7,0)
     .. controls (5.0,.48) and (6.3,-.50) .. (7.7,0);
  \fill[accent] (3.7,0) circle (1.5pt);
\end{tikzpicture}
\end{center}
\clearpage

\sectionline{1. Метод позиций (метод карточек)}
Когда требуется посчитать количество различных вариантов, разбейте каждый вариант на \textbf{последовательные позиции}: разряды числа, места в расписании, роли участников или участки пути.

В каждой позиции указывают \textbf{количество допустимых вариантов}, а затем перемножают эти количества.

\example{Сколько двузначных чисел можно составить из цифр $0,1,2,3$, если повторения разрешены?}
\[
\underbrace{3}_{\text{десятки}}\cdot
\underbrace{4}_{\text{единицы}}=\boxed{12}.
\]
Первая цифра не может быть нулём. Вторая может быть любой из четырёх цифр, в том числе уже использованной.

\sectionline{2. Правило умножения}
Если первый шаг можно выполнить $a$ способами и при \textbf{каждом} таком выборе второй шаг --- $b$ способами, общее число последовательностей равно $ab$. Для трёх шагов аналогично: $abc$.

\example{Между пунктами $A$ и $B$ есть 3 дороги, между $B$ и $C$ --- 4. Количество маршрутов $A\to B\to C$:}
\[3\cdot 4=\boxed{12}.\]
Слово «и» здесь означает: выбрать первую дорогу \textbf{и} вторую дорогу. Поэтому применяется умножение. Если варианты разбиты на непересекающиеся случаи «или», числа способов по случаям складывают.

\sectionline{3. Повторения и число претендентов}
Перед решением выясните, можно ли использовать один и тот же элемент несколько раз.
\begin{itemize}
\item \textbf{Повторения разрешены:} количество вариантов на следующем шаге сохраняется, если прочие условия его не ограничивают.
\item \textbf{Повторения запрещены:} использованные элементы исключаются. Если доступно $n$ разных элементов, то для трёх последовательных позиций получаем $n(n-1)(n-2)$.
\end{itemize}

\example{Сколько трёхзначных чисел без одинаковых цифр можно составить из $0,1,2,3$?}
\[
\underbrace{3}_{\text{сотни}}\cdot
\underbrace{3}_{\text{десятки}}\cdot
\underbrace{2}_{\text{единицы}}=\boxed{18}.
\]
После выбора первой цифры остаются три допустимые цифры: среди них может быть ноль.

\sectionline{4. Ограничения на отдельные позиции}
\keyidea{Сначала заполните самые ограниченные позиции: фиксированную первую или последнюю цифру, требование чётности, кратности, запрет на ноль. Затем считайте остальные позиции.}

\example{Из цифр $1,2,3,4,5$ составляют пятизначные числа без повторений, оканчивающиеся цифрой $3$. Последнюю позицию заполняем единственным способом, остальные --- $4!$ способами:}
\[4\cdot3\cdot2\cdot1\cdot1=\boxed{24}.\]

\textbf{Важно.} У числа старший разряд не может быть нулём. У цифрового \emph{кода} ведущие нули могут быть допустимы, если это прямо разрешено условием.

\sectionline{5. Факториал и перестановки}
Для натурального $n$ запись $n!$ означает произведение всех натуральных чисел от $n$ до 1:
\[
n!=n(n-1)(n-2)\cdots 2\cdot1,\qquad 0!=1.
\]
Если $n$ \textbf{различных} объектов размещают по $n$ позициям, используя каждый ровно один раз, вариантов $n!$.

\example{Пять разных школьных предметов нужно расставить по пяти урокам.}
\[5!=5\cdot4\cdot3\cdot2\cdot1=\boxed{120}.\]

\example{Из 18 человек выбирают физорга, культурного организатора и казначея. Один человек не занимает две должности.}
\[18\cdot17\cdot16=\boxed{4896}.\]
Здесь должности различаются: перестановка выбранных людей меняет результат.

\sectionline{6. Когда порядок имеет значение}
Сравните два действия:
\begin{itemize}
\item \textbf{Назначить} двух человек председателем и секретарём: порядок важен. Для $10$ человек это $10\cdot9=90$ способов.
\item \textbf{Выбрать} двоих из десяти в одну команду: порядок не важен. Каждый состав команды учтён дважды; получается $\dfrac{10\cdot9}{2!}=\boxed{45}$ способов.
\end{itemize}
Для выбора $k$ различных объектов из $n$ разных объектов \textbf{без учёта порядка} можно посчитать упорядоченные последовательности и разделить на $k!$:
\[
\frac{n(n-1)\cdots(n-k+1)}{k!},\qquad 1\le k\le n.
\]
Это рассуждение предполагает, что выбранные объекты различны и каждый набор допускает все $k!$ перестановок.

\sectionline{7. Если объекты должны стоять рядом}
\example{Среди восьми разных книг две образуют двухтомник и должны стоять рядом. Остальные шесть книг различны.}

Считаем два тома \textbf{одним блоком}. Тогда размещаем семь объектов: $7!$ вариантов. Внутри блока два тома можно переставить $2!$ способами:
\[7!\cdot2!=\boxed{10080}.\]
Нельзя забывать о внутреннем порядке томов, если он не закреплён условием.

\sectionline{8. Геометрический подсчёт: диагонали $n$-угольника}
Соединим любые две различные вершины. Упорядоченных пар $n(n-1)$, но отрезки $AB$ и $BA$ совпадают, поэтому разных отрезков $\frac{n(n-1)}2$.

Из них $n$ --- стороны многоугольника. Следовательно, число диагоналей:
\[
D=\frac{n(n-1)}2-n=\boxed{\frac{n(n-3)}2},\qquad n\ge3.
\]
\example{У шестиугольника $D=\frac{6\cdot3}{2}=9$, у двадцатиугольника $D=\frac{20\cdot17}{2}=170$.}

\sectionline{9. Делимость и ограничения на последние разряды}
Если число делится на $4$, число, образованное двумя его последними цифрами, делится на $4$. Для делимости на $5$ последняя цифра должна быть $0$ или $5$.

\example{Сколько шестизначных чисел без повторяющихся цифр, кратных $4$, можно составить из $1,3,5,6,7,9$?}

Возможные окончания: $16,36,56,76,96$ --- всего пять. Для каждого окончания остаются четыре цифры, которые можно расставить $4!$ способами:
\[5\cdot4!=\boxed{120}.\]

\clearpage
\thispagestyle{fancy}
\begin{center}{\Large\bfseries\color{darkaccent}Алгоритм: метод позиций}\end{center}
\vspace{1.0cm}
\begin{center}
\begin{tikzpicture}[
  >=Latex,
  line/.style={-Latex,draw=darkaccent,line width=.9pt},
  stagebox/.style={draw=accent!85!black,rounded corners=4pt,align=center,text width=9cm,minimum height=1.0cm,inner sep=6pt,font=\normalsize},
  decision/.style={diamond,draw=accent!85!black,aspect=3.2,text width=4.5cm,align=center,inner sep=2pt,font=\small},
  branch/.style={draw=accent!85!black,rounded corners=4pt,align=center,text width=4.3cm,minimum height=1.25cm,inner sep=4pt,font=\small},
  endnode/.style={draw=accent!85!black,ellipse,align=center,minimum width=3cm,minimum height=.8cm,font=\normalsize}
]
\node[endnode] (start) at (0,0) {Начало};
\node[stagebox] (positions) at (0,-1.6) {Выделите позиции и определите допустимые варианты для каждой};
\node[decision] (repeat) at (0,-4.1) {Повторения разрешены?};
\node[branch] (yes) at (-3.55,-6.7) {Число претендентов может сохраняться};
\node[branch] (no) at (3.55,-6.7) {Использованный элемент исключается};
\node[stagebox] (multiply) at (0,-9.1) {Учтите особые условия и перемножьте количество вариантов на позициях};
\node[decision] (order) at (0,-11.65) {Важен порядок?};
\node[branch] (orderyes) at (-3.55,-14.2) {Полученное число --- ответ};
\node[branch] (orderno) at (3.55,-14.2) {Для выбора $k$ разных объектов без повторений разделите на $k!$};
\node[endnode] (finish) at (0,-16.4) {Проверьте ответ};
\draw[line] (start) -- (positions);
\draw[line] (positions) -- (repeat);
\draw[line] (repeat.west) -| node[near start,above,sloped]{\scriptsize да} (yes.north);
\draw[line] (repeat.east) -| node[near start,above,sloped]{\scriptsize нет} (no.north);
\draw[draw=darkaccent,line width=.9pt] (yes.south) |- (0,-7.9);
\draw[draw=darkaccent,line width=.9pt] (no.south) |- (0,-7.9);
\draw[line] (0,-7.9) -- (multiply.north);
\draw[line] (multiply) -- (order);
\draw[line] (order.west) -| node[near start,above,sloped]{\scriptsize да} (orderyes.north);
\draw[line] (order.east) -| node[near start,above,sloped]{\scriptsize нет} (orderno.north);
\draw[draw=darkaccent,line width=.9pt] (orderyes.south) |- (0,-15.6);
\draw[draw=darkaccent,line width=.9pt] (orderno.south) |- (0,-15.6);
\draw[line] (0,-15.6) -- (finish.north);
\end{tikzpicture}
\end{center}
\vspace{1.0em}
\noindent\small\textcolor{softgray}{Примечание. Деление на $k!$ применимо при выборе разных объектов без повторений. Если повторы разрешены, а порядок не важен, нужен отдельный подсчёт.}
\clearpage

\begin{center}{\Large\bfseries\color{darkaccent}Тренировочный блок}\par\vspace{0.1cm}{\small Путь А · 10 заданий · от простого к сложному}\end{center}
\vspace{0.3cm}
В каждом задании запишите множители по позициям и краткое обоснование. Если порядок выбора не учитывается, объясните, на какое число делите и почему.

\subline{Базовый уровень}
\begin{enumerate}[label=\textbf{\arabic*.},itemsep=11pt]
\task{Сколько различных двузначных чисел можно составить из цифр $0,1,3,5,7$, если цифры могут повторяться?}
\task{Сколько четырёхзначных чисел без повторяющихся цифр можно составить из цифр $0,1,2,3,4$?}
\task{Из пункта $A$ в $B$ ведут четыре дороги, из $B$ в $C$ --- три, из $C$ в $D$ --- две. Сколько различных маршрутов из $A$ в $D$ через $B$ и $C$ существует?}
\task{Шесть разных учебных предметов нужно расставить в расписании из шести уроков по одному на урок. Сколько вариантов расписания возможно?}
\end{enumerate}
\vspace{0.4cm}
\subline{Уровень выше среднего}
\begin{enumerate}[label=\textbf{\arabic*.},start=5,itemsep=13pt]
\task{Из 12 участников выбирают председателя, заместителя и секретаря. Один участник не может занять две должности. Сколько способов назначения существует?}
\task{Сколько шестизначных чисел существует, если цифры на первом, третьем и пятом местах попарно различны, а остальные цифры могут повторяться? Используются цифры от $0$ до $9$.}
\task{Сколько пятизначных чисел, кратных $4$, можно составить из цифр $1,2,3,5,8$, используя каждую цифру ровно один раз?}
\task{На полке стоят семь разных книг. Две из них --- тома одного произведения; они должны стоять рядом. Сколько различных расстановок возможно, если порядок двух томов не ограничен?}
\end{enumerate}
\vspace{0.4cm}
\subline{Повышенный уровень}
\begin{enumerate}[label=\textbf{\arabic*.},start=9,itemsep=13pt]
\task{Из девяти различных книг нужно выбрать три для поездки. Порядок выбора не важен. Сколькими способами это можно сделать?}
\task{Сколько диагоналей у выпуклого пятнадцатиугольника? Выведите и примените формулу через общее число отрезков между вершинами.}
\end{enumerate}

\clearpage
\begin{center}{\Large\bfseries\color{darkaccent}Ответы к тренировочному блоку}\end{center}
\vspace{0.9cm}
\renewcommand{\arraystretch}{1.45}
\begin{tabularx}{\linewidth}{@{}p{1.2cm}X@{}}
\toprule
\textbf{№} & \textbf{Ответ и контрольное выражение}\\
\midrule
1 & $20;\quad 4\cdot5$.\\
2 & $96;\quad 4\cdot4\cdot3\cdot2$.\\
3 & $24;\quad 4\cdot3\cdot2$.\\
4 & $720;\quad 6!$.\\
5 & $1320;\quad 12\cdot11\cdot10$.\\
6 & $648000;\quad 9\cdot10\cdot9\cdot10\cdot8\cdot10$.\\
7 & $24;\quad 4\cdot3!$ (окончания $12,28,32,52$).\\
8 & $1440;\quad 6!\cdot2!$.\\
9 & $84;\quad \dfrac{9\cdot8\cdot7}{3!}$.\\
10 & $90;\quad \dfrac{15\cdot14}{2}-15=\dfrac{15\cdot12}{2}$.\\
\bottomrule
\end{tabularx}
\par\vspace{1.0cm}
\sectionline{Итоговая памятка}
\renewcommand{\arraystretch}{1.4}
\begin{tabularx}{\linewidth}{@{}p{.37\linewidth}X@{}}
\toprule
\textbf{Сигнал в условии} & \textbf{Как действовать}\\
\midrule
Несколько последовательных шагов & Перемножить количество вариантов для каждого шага.\\
Цифры не повторяются & После каждого выбора уменьшать число доступных цифр.\\
Первая цифра числа & Исключить ноль из первого разряда.\\
Все разные объекты переставляют & Использовать факториал.\\
Порядок выбора не важен & Устранить повторный учёт перестановок выбранных объектов.\\
Объекты должны стоять рядом & Объединить их в блок, затем учесть перестановки внутри блока.\\
Диагонали многоугольника & Посчитать пары вершин без порядка и вычесть стороны.\\
\bottomrule
\end{tabularx}
\par\vspace{0.7cm}
\noindent{\color{darkaccent}\bfseries Перед ответом проверьте:} повторения; важность порядка; ограничения на разряды; повторный учёт одного выбора.
\end{document}
